\cotefile{157-5}
\scanpage{69}\mspage{56}
% archive sheet number: 68

$f\in W(I^{\unsure{o|\circ}})$, alors $i\in W(I^o)$ donc $i\in W(I^o)\cap C_{I^o}=(TC)_{I^o}$.
Au total, \deleted{\unsure{tout}} $f\in W(I^o)_c$ est composé $pi$\inserted{ dans $M(I^o)_c$},
où $i\in(TC)_{I^o}$ et $p$ est inverse à gauche
d'un $j\in TC$. Comme $u_!$ transforme
$(TC)_{I^o}$ en $(TC)_{I'^o}$, il transforme $f$
en $f'=p'i'$, avec $i'\in(TC)_{I'^o}\subset W_{I'^o}$,
et $p'$ inverse à gauche d'un $j'=u_!(j)$
qui est \unsure{également} dans $W_{I'^o}$. Mais alors
$\deleted{\unsure{j'}}\unsure{p'|j'}\in W(I'^o)$, donc aussi $f'=p'i'\in W_{I'^o}$,
ce qui prouve que $u_!(W(I^o)_c)\subset W(I'^o)$.

La dernière assertion provient du \deleted{\unsure{fait}}\inserted{lemme}
suivant, appliqué à $(M(I^o),W(I^o),C_{I^o})$ \unsure{ald} $(M,W,C)$\unsure{.}

\uline{Lemme} \uline{3.4}. Soit $(M,W,C)$ une catégorie
à cofibrations. Alors l'inclusion $M_c\subset M$
induit une équivalence \deleted{\unsure{$\mathrm{Ho}_W$\illegible{1}}}
\[
M_c W_c^{-1}\overset{\approx}{\unsure{\to|\text{---}|\leftrightarrow}}MW^{-1}=\mathrm{Ho}_W\,.
\]

Nous donnerons la démonstration (par
\unsure{A la J|A le J|A et J|A 6 J|A \& J}) ultérieurement\unsure{,|.}

Cela achève la démonstration de
3.3.

3.5. Il se pose la question si \deleted{\unsure{le}} le foncteur
$u_!^W$ de (3.3.1) est bien un adjoint à g.
